\section{Burst Splitter}
\label{sec:splitter}
\subsection{Purpose and bounds}

The Burst Splitter converts one accepted AXI transaction into up to 64 MC
request packets.  It is shared by the read and write request paths.  On the
write path (\texttt{AWSIZE = 0..6}) it emits one packet per AXI write beat; on
the read path (\texttt{ARSIZE = 0..6}) it emits one packet per distinct
aligned 64-B line touched by the burst.  Each packet carries a 512-bit
payload on a 64-B byte-lane grid; there is no \texttt{beat\_cnt} field in the
packet.

The configured \texttt{row\_size} is a run-time parameter and is at least
1~KB.  AXI4 prohibits a burst from crossing a 4~KB boundary, so a read spans
at most 4~KB, i.e.\ at most
\[
 \texttt{MAX\_PKTS} = 4\,\mathrm{KB} / 64\,\mathrm{B} = 64
\]
64-B MC packets.  A write emits \texttt{AWLEN + 1} MC packets, bounded to the
same \texttt{MAX\_PKTS = 64} --- automatic for full-width writes and a
design-time assertion \texttt{AWLEN + 1 $\le$ 64} for narrow writes
(\S\ref{sec:overview}).  A transaction can still cross at most three 1~KB row
boundaries; row crossings affect each packet's address, not the packet
count.

\subsection{Interfaces and state}

At transaction acceptance, the selected read or write ROB allocates
\texttt{ROB\_INDEX[7:0]}.  The splitter accepts that index with the AXI
metadata and retains it for the whole transaction.  Stage~1 determines
\texttt{exp\_pkts}, the number of 64-B MC packets, and presents it to the
corresponding ROB; this is the sole source of the ROB's \texttt{exp\_pkts}.
On the write path there is one MC packet per AXI write beat regardless of
\texttt{AWSIZE} (\S\ref{sec:overview}), so \texttt{exp\_pkts = AWLEN + 1}.  On
the read path \texttt{exp\_pkts} is the number of distinct 64-B lines the
burst touches, computed per burst type in Stage~1; it is \emph{not}
\texttt{ARLEN + 1}.

\begin{center}
\begin{tabular}{lll}
\toprule
\textbf{State} & \textbf{Width/depth} & \textbf{Meaning}\\
\midrule
\texttt{ROB\_INDEX} & 8 bits & Transaction index retained while packets issue\\
\texttt{exp\_pkts} & 7 bits & Number of 64-B MC packets for the transaction, \texttt{1..64}\\
\texttt{pkt\_idx} & 6 bits & \texttt{PACKET\_NUMBER} of the packet currently being emitted, \texttt{0..exp\_pkts-1} (write: AXI beat index; read: 64-B line index)\\
\bottomrule
\end{tabular}
\end{center}

\noindent The internal structure Stage~1 uses to record the per-packet
address layout (a per-beat address generator, and how row boundaries are
tracked within it) is not fixed by this document.

\subsection{Stage 1 --- packet-layout construction}

Stage~1 receives address, burst length, size, burst type, AXID, packet type,
and \texttt{ROB\_INDEX}.  It computes \texttt{exp\_pkts} --- the number of MC
packets the transaction requires --- and the per-packet address layout.  The
per-transaction packet count is bounded at 64 (\S\ref{sec:overview}).

\textbf{Write path} (\texttt{AWSIZE = 0..6}): one MC packet per AXI write
beat, \texttt{exp\_pkts = AWLEN + 1}, \texttt{PACKET\_NUMBER} $=$ the AXI beat
index.  Packet \texttt{N} carries \texttt{WDATA[511:0]} and \texttt{WMASK[63:0]}
verbatim from that beat's \texttt{WSTRB}; its address is beat \texttt{N}'s AXI
address (\texttt{INCR}: \texttt{base + N*(1 << AWSIZE)}; \texttt{WRAP}: the
wrap-table entry; \texttt{FIXED}: \texttt{base}), mapped to its containing
64-B line for \texttt{daddr}.  Narrow beats are not coalesced --- two beats in
the same 64-B line are still two MC packets with partial \texttt{WMASK}.

\textbf{Read path} (\texttt{ARSIZE = 0..6}): one packet per distinct 64-B
line touched.  With \texttt{NB = 1 << ARSIZE}:
\begin{itemize}[noitemsep]
  \item \texttt{INCR}: \texttt{lo = ARADDR};
        \texttt{hi = (ARADDR \& \textasciitilde(NB - 1)) + (ARLEN + 1)*NB - 1};
        \texttt{first\_line = lo >> 6};
        \texttt{exp\_pkts = (hi >> 6) - first\_line + 1}.
  \item \texttt{WRAP}: \texttt{Total = (ARLEN + 1) << ARSIZE};
        \texttt{WrapBoundary = ARADDR \& \textasciitilde(Total - 1)};
        \texttt{first\_line = WrapBoundary >> 6};
        \texttt{exp\_pkts = (Total <= 64) ? 1 : (Total >> 6)}.
  \item \texttt{FIXED}: \texttt{lo = ARADDR};
        \texttt{first\_line = lo >> 6}; \texttt{exp\_pkts = 1}.
\end{itemize}
Read packet \texttt{N} (\texttt{PACKET\_NUMBER = N}, the line index)
addresses the 64-B line \texttt{first\_line + N}.  Unaligned first transfers
are handled per AXI4 semantics by the Read Serializer
(\S\ref{sec:read-serializer}): they change which bytes the first beat
carries, not the line set, which \texttt{first\_line}/\texttt{exp\_pkts}
already captures.

For a read allocation, Stage~1 provides \texttt{exp\_pkts} to the read ROB
before its \texttt{ROB\_INDEX} becomes visible at the relevant per-AXID Slot
FIFO.  The write path follows the identical ordering.

\subsection{Stage 2 --- sequential emission}

Stage~2 emits the transaction's packets in increasing \texttt{pkt\_idx}
order.  For each accepted packet it forms
\[
 \texttt{GLOBAL\_TAG[13:0]=\{ROB\_INDEX[7:0],PACKET\_NUMBER[5:0]\}},
 \qquad \texttt{PACKET\_NUMBER=pkt\_idx}.
\]
Packet Formation receives the packet's 64-B line \texttt{addr}, packet
type, and \texttt{GLOBAL\_TAG}; it serializes one MC request packet, carrying
\texttt{WDATA[511:0]} + \texttt{WMASK[63:0]} for one AXI write beat (write) or
no payload (read).  No per-packet Read
Data Buffer handle is recorded --- the read ROB's contiguous
\texttt{rd\_dbuf\_base} range was allocated at ROB initialization and packet
\texttt{pkt\_idx}'s data lands at \texttt{rd\_dbuf\_base + pkt\_idx}.

Stage~2 advances \texttt{pkt\_idx} only on downstream packet acceptance.  It
stalls Stage~1 for the duration of a multi-packet transaction, preserving the
packet layout and its associated \texttt{ROB\_INDEX}.  After the final
accepted packet it returns to idle and permits the next Stage~1 transaction.

\subsection{WRAP and FIXED addressing}

This subsection describes the \emph{read}-path per-line layout.  On the write
path there is one MC packet per AXI write beat for all three burst types, and
packet \texttt{N}'s address is beat \texttt{N}'s AXI address (Stage~1, Write
path).

For a read \texttt{INCR} transaction, packet \texttt{N}'s 64-B line address is
\texttt{first\_line + N} lines.  For a \texttt{WRAP} transaction the line set
is the aligned \texttt{Total}-byte window based at \texttt{WrapBoundary}
(Stage~1); packet \texttt{N} addresses line \texttt{first\_line + N} within
that window, so the MC packets still cover a contiguous line range.
\texttt{FIXED} produces a single line (\texttt{exp\_pkts = 1}).  The
per-line \texttt{PACKET\_NUMBER} assignment (\texttt{0..exp\_pkts-1}) is
identical across burst types; only \texttt{first\_line}/\texttt{exp\_pkts}
differ.  Reconstructing the original AXI beat addresses --- \texttt{WRAP}
wrap-around and narrow strides included --- is the Read Serializer's
responsibility (\S\ref{sec:read-serializer}), not the splitter's.

\subsection{Backpressure and error packets}

\begin{itemize}[noitemsep]
  \item A full selected ROB Slot FIFO, an empty selected ROB Free List, or
        unavailable read Data Buffer capacity prevents acceptance upstream as
        appropriate to that path.
  \item Packet-FIFO backpressure holds the current packet and its
        \texttt{pkt\_idx}; it does not alter \texttt{exp\_pkts}.
  \item Error read and error write transactions use the same per-packet
        \texttt{GLOBAL\_TAG} mechanism.  Their packet type is carried in
        every generated packet.
\end{itemize}

\subsection{Timing notes}

\begin{itemize}[noitemsep]
  \item \texttt{exp\_pkts} is the number of 64-B MC packets generated by this
        one transaction.  Write: the AXI write-beat count
        (\texttt{AWLEN + 1}).  Read: the count of distinct 64-B lines touched
        (Stage~1), generally less than \texttt{ARLEN + 1} for narrow bursts.
  \item \texttt{PACKET\_NUMBER} (\texttt{= pkt\_idx}) is assigned at packet
        emission, not allocation.
\end{itemize}
